% =====================================================================================
% Figure 1 -- wireframe authoring against the Gaussian-splat twin. Three panels, so the
% figure covers BOTH authoring techniques the section describes: scaled primitives for
% straight edges, and extruded splines for curves and trees.
%
% Drop-in replacement for the "TODO: insert Figure 1" block in
% \subsection{Capturing and Rendering the Digital Twin}.
%
% Requires in the preamble (acmart does not load it by default):
%     \usepackage{subcaption}
%
% Expected image files, relative to the main .tex:
%     figures/wireframe-authoring-exterior.jpg   (brick building, three-quarter view)
%     figures/wireframe-authoring-steps.jpg      (steps and columns, eye level)
%     figures/wireframe-authoring-tree.jpg       (tree limbs traced as splines)
%
% The three sources have different aspect ratios (roughly 4:3, 3:2, 5:3). Side by side at
% equal width they will render at unequal heights. Crop all three to a common ratio before
% exporting -- 3:2 loses the least from each.
%
% Reference in text as Figure~\ref{fig:wireframe}. Suggested spot: end of the paragraph
% beginning "With the Gaussian reconstruction loaded as a live tracing target...".
% A single-panel fallback is commented at the bottom of this file.
% =====================================================================================

\begin{figure*}[t]
  \centering
  \begin{subfigure}[b]{0.32\textwidth}
    \includegraphics[width=\linewidth]{figures/wireframe-authoring-exterior}
    \caption{Building exterior.}
    \label{fig:wireframe-exterior}
  \end{subfigure}
  \hfill
  \begin{subfigure}[b]{0.32\textwidth}
    \includegraphics[width=\linewidth]{figures/wireframe-authoring-steps}
    \caption{Steps and columns.}
    \label{fig:wireframe-steps}
  \end{subfigure}
  \hfill
  \begin{subfigure}[b]{0.32\textwidth}
    \includegraphics[width=\linewidth]{figures/wireframe-authoring-tree}
    \caption{Tree limbs.}
    \label{fig:wireframe-tree}
  \end{subfigure}
  \caption{Authoring views: the blue wireframe traced by hand against the photoreal
    Gaussian-splat twin loaded as a tracing target in the Unity scene view. The splat is a
    development artifact and is destroyed in non-Editor builds; the participant sees only
    the blue lines. (\subref{fig:wireframe-exterior}) Straight features are built from
    Unity primitive cylinders and capsules, scaled thin along two axes so each reads as a
    line: roof line, wall corners, window frames, and railings.
    (\subref{fig:wireframe-steps}) The same technique at close range on steps, low walls,
    column edges, handrails, and door frames. (\subref{fig:wireframe-tree}) Curved features
    have no matching primitive and are drawn as splines swept into tube mesh; the blue
    loops over the canopy are the traced trunk and limbs, and the curb crossing the
    foreground is a further spline. Throughout, only edges that carry the structure of the
    space are authored, and broad areas of brick, pavement, and foliage are left untraced.}
  \Description{Three screenshots from a Unity scene view, each showing a photorealistic
    three-dimensional scan of a campus site overlaid with thin blue lines of uniform color
    and brightness. The first shows a red brick building from a three-quarter angle, with
    blue lines along the roof line, the vertical corners of the walls, the frames of tall
    narrow windows, a metal railing, and the curved edge of the sidewalk. The second, taken
    at eye level, shows a flight of concrete steps rising to a covered entrance, with blue
    lines along each step edge, the low walls flanking the steps, the square columns
    supporting the roof, the metal handrails, and the frames of the glass doors and windows
    behind. The third shows a large leafy tree in a planting bed in front of a brick
    building, with long curving blue lines tracing the trunk and the spreading limbs through
    the canopy, and a further blue line following the raised concrete curb that separates
    the planting bed from a lawn in the foreground. In all three images, large areas of
    brick, pavement, grass, and foliage carry no lines at all.}
  \label{fig:wireframe}
\end{figure*}

% -------------------------------------------------------------------------------------
% Single-panel fallback, if three panels crowd the layout. Delete the block above.
% -------------------------------------------------------------------------------------
%
% \begin{figure}[t]
%   \centering
%   \includegraphics[width=\linewidth]{figures/wireframe-authoring-exterior}
%   \caption{Authoring view: the blue wireframe traced by hand against the photoreal
%     Gaussian-splat twin loaded as a tracing target in the Unity scene view. The splat is a
%     development artifact and is destroyed in non-Editor builds; the participant sees only
%     the blue lines. Only edges that carry the structure of the space are authored: the
%     roof line, wall corners, window frames, and railings, with the curved sidewalk curb
%     drawn as an extruded spline. Broad areas of brick and pavement are left untraced.}
%   \Description{A screenshot from a Unity scene view showing a photorealistic
%     three-dimensional scan of a red brick building, seen from a three-quarter angle,
%     overlaid with thin blue lines of uniform color and brightness. The lines follow the
%     roof line, the vertical corners of the walls, the frames of tall narrow windows, a
%     metal railing beside a planted area, and the curved edge of the sidewalk in the
%     foreground. Large areas of brick wall and pavement carry no lines at all.}
%   \label{fig:wireframe}
% \end{figure}
